\section{Erd\H{o}s Problem \#1154: a zero-dimensional field construction and an analytic-generator obstruction}

\subsection*{1) Formal restatement and status}
Does every $\alpha\in[0,1]$ occur as the Hausdorff dimension of a subring, or a subfield, of $\mathbb R$ in ZFC? The live page, checked 23 September 2026, marks the unrestricted question NOT DISPROVABLE: Mauldin gives the requested subfields assuming the continuum hypothesis. We do not assume CH here. The endpoint dimensions are immediate from $\mathbb Q$ and $\mathbb R$; the unresolved target is arbitrary $0<\alpha<1$ without that additional axiom.

We use one external theorem explicitly: Edgar--Miller prove that an analytic subring of $\mathbb R$ either has Hausdorff dimension zero or equals $\mathbb R$. This is a restriction on analytic rings, not a proof against arbitrary nonanalytic rings.

\subsection*{2) Objective}
A natural construction is to start with a thin Cantor set and take its generated field. We prove that extremely thin generators do give a proper uncountable field of dimension zero. We then show, using the stated external theorem, why positive-dimensional analytic generators cannot interpolate this construction to an intermediate dimension.

\subsection*{3) A self-contained uncountable zero-dimensional example}
Set
\[
 C=\left\{\sum_{n=2}^{\infty}\epsilon_n2^{-n!}:\epsilon_n\in\{0,1\}\right\}.
\]
The binary choices give a compact uncountable set. Indeed, for the first index where two choices differ, the leading summand exceeds the sum of all subsequent summands, so the two real numbers are distinct. Compactness follows from the uniformly convergent sum on the compact product $\{0,1\}^{\mathbb N}$.

\textbf{Lemma.} $C$ has upper box dimension zero, and so does $C^d$ for every finite $d$.

\emph{Proof.} For $k$ large and
\[
 2^{-(k+1)!}\leq\delta<2^{-k!},
\]
fixing the digits through index $k+1$ leaves a tail of length at most $2\,2^{-(k+2)!}$, smaller than $\delta$. Thus $C$ is covered by at most $2^{k+1}$ intervals of length $\delta$, up to an absolute harmless factor. Therefore
\[
 \frac{\log N(C,\delta)}{-\log\delta}\leq\frac{(k+1)\log2+O(1)}{k!\log2}\longrightarrow0.
\]
Taking products gives at most $2^{d(k+1)}$ cubes of side $\delta$, with diameter $\sqrt d\delta$. The same limit is zero. In particular each finite product has Hausdorff dimension zero. $\square$

Let $F=\mathbb Q(C)$ be the field generated by $C$. Every member is a value
\[
 P(x_1,\ldots,x_d)/Q(x_1,\ldots,x_d),\qquad x_i\in C,
\]
for some finite $d$, rational-coefficient polynomials $P,Q$, and a nonzero denominator. There are countably many such choices of $d,P,Q$.

For each positive integer $j$, restrict the domain to the compact set
\[
 K_{P,Q,j}=\{x\in C^d:|Q(x)|\geq1/j\}.
\]
On this set the quotient map is Lipschitz. To verify this without a convexity assumption on the restricted domain, write
\[
 \left|\frac{P(x)}{Q(x)}-\frac{P(y)}{Q(y)}\right|
 \leq j|P(x)-P(y)|+j^2|P(y)|\,|Q(x)-Q(y)|.
\]
The polynomial differences are Lipschitz and $P$ is bounded on a fixed box containing $C^d$. Hence the image of $K_{P,Q,j}$ is compact and has Hausdorff dimension zero. Every nonzero denominator is covered for some $j$, so $F$ is a countable union of those compact zero-dimensional images. Countable stability gives $\dim_H F=0$.

Thus $F$ is an uncountable proper subfield of $\mathbb R$, and even an $F_\sigma$ set. Uncountability follows from $C\subset F$; properness follows from its dimension being zero while $\dim_H\mathbb R=1$. All claims in this construction have been proved here.

\subsection*{4) Why analytic generators cannot give intermediate dimensions}
\textbf{Proposition, using Edgar--Miller.} If an analytic set $S\subseteq\mathbb R$ has positive Hausdorff dimension, its generated unital ring is $\mathbb R$, and consequently $\mathbb Q(S)=\mathbb R$.

\emph{Proof.} The generated unital ring is the union, over all finite $d$ and integer polynomials $P$, of the sets $P(S^d)$. Finite products and continuous images of analytic sets are analytic, and countable unions of analytic sets are analytic. Thus the generated ring is analytic. Its dimension is positive because it contains $S$. Edgar--Miller therefore force it to equal $\mathbb R$. The generated field contains that ring. $\square$

\textbf{Necessary condition for a hypothetical solution.} Any subring $R$ with $0<\dim_H R<1$ contains no positive-dimensional analytic subset. If $R$ is unital, such a subset would generate all of $\mathbb R$ inside $R$. Without a unital convention, its generated unital ring lies inside $R+\mathbb Z$. This is a ring and a countable union of translates of $R$, hence has the same Hausdorff dimension as $R$; it cannot contain all of $\mathbb R$ either. In particular every compact subset of $R$ must have dimension zero. This is a strong obstruction to a compact-fractal construction, but arbitrary nonanalytic sets need not have their dimension witnessed by compact subsets. Treating that last assertion as automatic would incorrectly solve the problem.

\subsection*{5) Adversarial verification}
The dimension estimate for $C$ controls all sufficiently small scales, not just the factorial subsequence. Hausdorff dimension alone would not generally give an upper bound for dimensions of products; we use upper box dimension zero to justify that step. Rational maps may have poles, so the proof restricts denominators away from zero and then takes a countable union. Finally, the regularity obstruction invokes precisely an analytic-ring theorem and does not extend it to arbitrary rings by assumption.

\subsection*{6) Final}
\textbf{UNRESOLVED.} We construct an uncountable zero-dimensional field in ZFC and isolate a rigorous obstruction to the most direct interpolation route. We have neither constructed an intermediate-dimensional field without CH nor proved independence of the requested existence statement. The missing ingredient must go beyond a positive-dimensional analytic generating set.

\subsection*{7) References}
\begin{itemize}
\item Current problem statement and status: \url{https://www.erdosproblems.com/1154}.
\item G. A. Edgar and C. Miller, \emph{Borel subrings of the reals}, Proc. Amer. Math. Soc. 131 (2003), 1121--1129; author-hosted text: \url{https://math.osu.edu/~edgar/preprints/subring/subringpre.pdf}.
\item R. D. Mauldin, \emph{Subfields of R with arbitrary Hausdorff dimension}, Math. Proc. Cambridge Philos. Soc. 161 (2016), 157--165. The CH theorem is context, not used in our construction.
\end{itemize}
No novelty over existing fractal-field constructions is claimed. This is an independently justified partial attempt. The completion percentage is a subjective assessment, not a theorem or a correctness probability.
\subsection*{8) COMPLETION ESTIMATE}
\textbf{COMPLETION: 10\%}
